% !TeX root = ../../Book3.tex
\section{Jacobi 三重积公式}
\label{b3:sec:2402}

级数把不同整数指标的贡献相加，适合研究变换和表示数；乘积把零点逐个列出，适合判断除法是否有意义。这里要把上一节的 theta 级数改写成乘积，并说明为什么无穷多个非零因子不会在极限中另外制造零点。

证明从有限多项式开始。先把有限乘积的系数算清楚，再把指标居中，使指数变成平方，最后在紧集上取极限。这条路线同时给出恒等式中的常数，避免只凭零点相同就猜测两个函数相等。

\subsection{三重积恒等式}
\label{b3:subsec:240201}

下面的有限恒等式起着普通二项式公式的作用。记号中的下标 \(Q\) 表示因子按等比数列变化；这里的系数一般是复数，不借用计数或正性来证明公式。

\Needspace{8\baselineskip}
\begin{lemma}[有限二项式恒等式]\label{b3:lem:24-finite-binomial}
设 \(Q\in\mathbb C\)、\(0<|Q|<1\)，并令
\[
 (Q;Q)_0=1,\qquad (Q;Q)_N=\prod_{j=1}^N(1-Q^j)
 \quad(N\ge1).
\]
对整数 \(N\ge0\) 及 \(0\le r\le N\)，定义
\[
 {N\brack r}_Q
 =\frac{(Q;Q)_N}{(Q;Q)_r(Q;Q)_{N-r}},
\]
并在 \(r<0\) 或 \(r>N\) 时规定 \({N\brack r}_Q=0\)。则对任意 \(t\in\mathbb C\)，
\begin{equation}\label{b3:eq:24-finite-binomial}
 \prod_{j=0}^{N-1}(1+tQ^j)
 =\sum_{r=0}^N Q^{r(r-1)/2}{N\brack r}_Q\,t^r.
\end{equation}
当 \(N=0\) 时，左边的空乘积和右边都等于一。
\end{lemma}
\begin{proof}
因 \(|Q|<1\)，所有分母因子都不为零。先证明递推关系
\[
 {N+1\brack r}_Q
 ={N\brack r}_Q+Q^{N+1-r}{N\brack r-1}_Q.
\]
在 \(1\le r\le N\) 时，把右边通分为分母
\((Q;Q)_r(Q;Q)_{N+1-r}\)，分子为
\[
 (Q;Q)_N\bigl((1-Q^{N+1-r})
              +Q^{N+1-r}(1-Q^r)\bigr)
 =(Q;Q)_{N+1}.
\]
在 \(r=0,N+1\) 时，越界系数为零的约定直接给出同一递推。

现在对 \(N\) 归纳。把 \(N\) 个因子的展开乘以 \(1+tQ^N\)，所得 \(t^r\) 的系数为
\[
 Q^{r(r-1)/2}{N\brack r}_Q
 +Q^N Q^{(r-1)(r-2)/2}{N\brack r-1}_Q.
\]
提出 \(Q^{r(r-1)/2}\) 后，括号中的式子就是上述递推的右边，因此得到 \(N+1\) 的公式。
\end{proof}

例如 \(N=2\) 时，公式是
\[
 (1+t)(1+tQ)=1+(1+Q)t+Qt^2.
\]
额外的因子 \(Q^{r(r-1)/2}\) 不能省去：选择两个因子中的 \(t\) 项时，所得系数是 \(Q\)，而不是一。

将有限公式居中后，就得到下面的\term[jacobi-sanchongji]{Jacobi 三重积公式}[Jacobi triple product identity]。

\Needspace{8\baselineskip}
\begin{theorem}[Jacobi 三重积]\label{b3:thm:24-triple-product}
设 \(q,w\in\mathbb C\)，且 \(0<|q|<1\)、\(w\ne0\)。则
\[
 \sum_{n\in\mathbb Z}q^{n^2}w^n
 =\prod_{m=1}^{\infty}
   (1-q^{2m})(1+wq^{2m-1})(1+w^{-1}q^{2m-1}).
\]
固定 \(q\) 后，级数绝对正规收敛，乘积正规收敛，二者都在
\(\mathbb C\setminus\{0\}\) 上定义关于 \(w\) 的全纯函数。
\end{theorem}
\begin{proof}
在式\eqref{b3:eq:24-finite-binomial}中，把因子个数取为 \(2N\)，再取
\[
 Q=q^2,\qquad t=wq^{-(2N-1)}.
\]
左边各因子中 \(q\) 的指数依次为
\(-(2N-1),-(2N-3),\ldots,2N-1\)。对负指数的那一半使用
\[
 1+wq^{-(2j-1)}
 =wq^{-(2j-1)}(1+w^{-1}q^{2j-1}).
\]
提出的全部单项式相乘为 \(w^Nq^{-N^2}\)，因为前 \(N\) 个正奇数的和为 \(N^2\)。

右边第 \(\ell\) 项含有
\[
 q^{\ell(\ell-1)}w^\ell q^{-(2N-1)\ell}
 =w^\ell q^{\ell^2-2N\ell}.
\]
因此两边乘以 \(w^{-N}q^{N^2}\)，并令 \(r=\ell-N\)，就得到
\begin{equation}\label{b3:eq:24-centered-product}
 \prod_{j=1}^N(1+wq^{2j-1})(1+w^{-1}q^{2j-1})
 =\sum_{r=-N}^N q^{r^2}w^r{2N\brack N+r}_{q^2}.
\end{equation}
平方指数来自 \(\ell^2-2N\ell+N^2=(\ell-N)^2\)，而双边指标来自把长度为 \(2N\) 的有限区间移到原点两侧。

接着说明极限为什么可以进入右边的和。由
\(\sum_{j\ge1}|q|^{2j}<\infty\) 和\cref{b3:prop:20-absolute-product}的非零乘积判据，
\[
 A=\prod_{j=1}^{\infty}(1-q^{2j})
\]
存在且非零。对每个固定整数 \(r\)，当 \(N\to\infty\) 时，商式中的三个有限乘积都趋于 \(A\)，故
\[
 {2N\brack N+r}_{q^2}\longrightarrow A^{-1}.
\]
这里还需要一个不随 \(N,r\) 变化的界。记
\[
 d=\prod_{j=1}^{\infty}(1-|q|^{2j})>0,\qquad
 M=\prod_{j=1}^{\infty}(1+|q|^{2j})<\infty.
\]
任一分母有限乘积的模不小于 \(d\)，分子有限乘积的模不大于 \(M\)，所以
\[
 \left|{2N\brack N+r}_{q^2}\right|
 \le Md^{-2}\qquad(-N\le r\le N).
\]

固定紧圆环 \(a\le|w|\le b\)，其中 \(0<a\le b<\infty\)，令
\(B=\max\{b,a^{-1},1\}\)。把区间外的系数补为零后，右边每一项的模都不超过
\[
 Md^{-2}|q|^{r^2}B^{|r|},\qquad r\in\mathbb Z.
\]
这个双边数列可和，因为平方衰减压过固定的指数增长。于是先把尾和统一压小，再对剩余有限项取极限，得到右边在紧圆环上一致趋于
\(A^{-1}\sum_{r\in\mathbb Z}q^{r^2}w^r\)。

左边也在紧圆环上一致收敛：两个因子与一的差分别被
\(b|q|^{2j-1}\) 和 \(a^{-1}|q|^{2j-1}\) 控制，控制数列均可和。
\cref{b3:prop:20-normal-product}因而适用。对式\eqref{b3:eq:24-centered-product}取极限，再乘以 \(A\)，即得所需公式及收敛结论。
\end{proof}

固定 \(r\) 的系数极限并不足以独自证明定理，因为求和区间还在扩大。证明中的共同上界专门控制那些随 \(N\) 向两端移动的项。另一方面，三重积的某个因子可以为零；正规收敛所保证的是充分远的尾部非零，并没有排除有限头部给出的零点。

\subsection{乘积展开}
\label{b3:subsec:240202}

现在恢复 \(\tau\in\mathfrak H\)、\(q=e^{\pi\mathrm i\tau}\)。本章的空间变量是 \(z\)，因此周期中出现 \(\pi\)，而不是一。以下所有
\(q^{1/4}\) 都表示 \(e^{\pi\mathrm i\tau/4}\)；把 \(q\) 单独看成穿孔圆盘中的变量时，不能另选一个全局的四次根来替换它。

\Needspace{8\baselineskip}
\begin{proposition}[四种 theta 函数的乘积]\label{b3:prop:24-theta-products}
设 \(\tau\in\mathfrak H\)、\(q=e^{\pi\mathrm i\tau}\)，并以
\(q^{1/4}=e^{\pi\mathrm i\tau/4}\) 规定分数幂。则上一节定义的四种 theta 函数对每个 \(z\in\mathbb C\) 满足
\begin{align}
 \theta_1(z,\tau)
 &=2q^{1/4}\sin z
   \prod_{n\ge1}(1-q^{2n})
   (1-2q^{2n}\cos2z+q^{4n}),
   \label{b3:eq:24-theta1-product}\\
 \theta_2(z,\tau)
 &=2q^{1/4}\cos z
   \prod_{n\ge1}(1-q^{2n})
   (1+2q^{2n}\cos2z+q^{4n}),
   \label{b3:eq:24-theta2-product}\\
 \theta_3(z,\tau)
 &=\prod_{n\ge1}(1-q^{2n})
   (1+2q^{2n-1}\cos2z+q^{4n-2}),
   \label{b3:eq:24-theta3-product}\\
 \theta_4(z,\tau)
 &=\prod_{n\ge1}(1-q^{2n})
   (1-2q^{2n-1}\cos2z+q^{4n-2}).
   \label{b3:eq:24-theta4-product}
\end{align}
固定 \(\tau\) 后，乘积关于 \(z\) 正规收敛，并可对 \(z\) 作任意固定阶数的逐项微分。其收敛在 \((z,\tau)\) 的紧集上也一致。
\end{proposition}
\begin{proof}
在三重积中分别取 \(w=e^{2\mathrm iz}\) 和 \(w=-e^{2\mathrm iz}\)，得到 \(\theta_3,\theta_4\) 的式子。其中两个含 \(w\) 的因子相乘，分别给出正号和负号的余弦项。

对半整数指标，先在级数中提取共同因子：
\[
 \begin{aligned}
 \theta_2(z,\tau)
 &=q^{1/4}e^{\mathrm iz}
   \sum_{n\in\mathbb Z}q^{n^2}(qe^{2\mathrm iz})^n,\\
 \theta_1(z,\tau)
 &=-\mathrm i q^{1/4}e^{\mathrm iz}
   \sum_{n\in\mathbb Z}q^{n^2}(-qe^{2\mathrm iz})^n.
 \end{aligned}
\]
第一行的两个指数因子成为
\(1+q^{2m}e^{2\mathrm iz}\) 和
\(1+q^{2m-2}e^{-2\mathrm iz}\)。
后一个乘积在 \(m=1\) 时提供 \(1+e^{-2\mathrm iz}\)；将剩余指标减一，便与第一个乘积配对。外面的因子化为
\[
 e^{\mathrm iz}(1+e^{-2\mathrm iz})=2\cos z.
\]
第二行同样给出负号的两个指数因子，而外面的因子为
\[
 -\mathrm i e^{\mathrm iz}(1-e^{-2\mathrm iz})=2\sin z.
\]
这就解释了四个公式中不同的首因子及符号。

若 \((z,\tau)\) 限制在一个紧集，则有 \(|q|\le r<1\) 及
\(|e^{\pm2\mathrm iz}|\le C\)。把每个二次因子拆为
\(1\pm q^m e^{2\mathrm iz}\) 与 \(1\pm q^m e^{-2\mathrm iz}\)，其偏离一的量被 \(Cr^m\) 控制。故乘积的正规收敛及有限分组均由\cref{b3:prop:20-normal-product}保证。固定 \(\tau\) 后，在稍大的 \(z\) 紧集上使用同一估计及 Cauchy 公式，就得到各阶导数的局部一致收敛。
\end{proof}

乘积表示给出的不只是一些显然存在的零点，还包括所有零点及其重数。后文把 theta 函数放在分母时，将用到这一完整信息。

\Needspace{8\baselineskip}
\begin{proposition}[theta 零点及常数非零性]\label{b3:prop:24-theta-zeros}
设 \(\tau\in\mathfrak H\)，令
\(\Lambda=\pi\mathbb Z+\pi\tau\mathbb Z\)。四个整函数
\(z\mapsto\theta_j(z,\tau)\) 的零点集分别为
\[
 \begin{aligned}
 Z(\theta_1)&=\Lambda,
 &Z(\theta_2)&=\frac{\pi}{2}+\Lambda,\\
 Z(\theta_3)&=\frac{\pi+\pi\tau}{2}+\Lambda,
 &Z(\theta_4)&=\frac{\pi\tau}{2}+\Lambda.
 \end{aligned}
\]
其中每一个零点都是单零点。特别地，
\(\theta_2(0,\tau)\)、\(\theta_3(0,\tau)\)、\(\theta_4(0,\tau)\)
均不为零，而 \(\theta_1(0,\tau)=0\) 且
\(\partial_z\theta_1(0,\tau)\ne0\)。
\end{proposition}
\begin{proof}
在 \(\theta_1\) 的乘积中，因子 \(1-q^{2n}\) 从不为零。
首因子 \(\sin z\) 的零点是 \(\pi\mathbb Z\)，且都单重。另外两个因子为
\[
 1-q^{2n}e^{2\mathrm iz},\qquad
 1-q^{2n}e^{-2\mathrm iz}.
\]
它们的零点分别是 \(-n\pi\tau+\pi\mathbb Z\) 和
\(n\pi\tau+\pi\mathbb Z\)。在各自的零点处，导数为
\(-2\mathrm i\) 或 \(2\mathrm i\)，所以也都是单零点。
因 \(\operatorname{Im}\tau>0\)，这些水平直线彼此不同，且不与
\(\sin z\) 的零点重叠。于是每个格点恰好来自一个单零因子。

对 \(\theta_2\)，首因子换成 \(\cos z\)，两个指数因子换成正号，零点的实部相应平移 \(\pi/2\)，得到 \(\pi/2+\Lambda\)。
对 \(\theta_3,\theta_4\)，指数为 \(2n-1\)，所以虚方向的平移为
\(\pm(n-\tfrac12)\pi\tau\)。正号因子再增加 \(\pi/2\) 的水平平移，负号因子不增加，分别得到所列的两个集合。各因子的导数仍非零，不同因子的零点仍不重叠。

最后，正规乘积在每个小邻域都可写成有限头部乘一个非零的全纯尾部；这是\cref{b3:prop:20-product-zero-orders}的具体适用情形。因此没有遗漏的额外零点，也没有因无限乘积而增加的重数。原点不属于后三个格的平移：比较虚部先确定 \(\tau\) 的系数，再比较实部即可验证。常数非零性随之成立。
\end{proof}

四个零点集在一个格胞中分别占据四个半周期位置。准周期平移会改变函数值的非零乘子，但不会改变零点的重数；因此这里描述的是整张平面上的零点排列，而不是只在一个基本区域内的局部结论。

\subsection{特殊值}
\label{b3:subsec:240203}

令 \(z=0\)。为减轻记号，以下把 \(\theta_j(0,\tau)\) 简写成
\(\theta_j\)，但撇号始终表示对空间变量 \(z\) 求导。由乘积公式，
\[
 \begin{aligned}
 \theta_2&=2q^{1/4}
       \prod_{n\ge1}(1-q^{2n})(1+q^{2n})^2,\\
 \theta_3&=\prod_{n\ge1}(1-q^{2n})(1+q^{2n-1})^2,\\
 \theta_4&=\prod_{n\ge1}(1-q^{2n})(1-q^{2n-1})^2.
 \end{aligned}
\]
这些是关于参数 \(\tau\) 的函数；“常数”只表示空间变量已经取零。

三个非零常数还决定奇函数在原点的导数，这一关系称为\term[jacobi-daoshu-hengdengshi]{Jacobi 导数恒等式}[Jacobi derivative identity]。

\Needspace{6\baselineskip}
\begin{proposition}[Jacobi 导数恒等式]\label{b3:prop:24-jacobi-derivative}
设 \(\tau\in\mathfrak H\)，四个 theta 函数采用本章的级数归一化，则
\begin{equation}\label{b3:eq:24-derivative-identity}
 \partial_z\theta_1(0,\tau)
 =\theta_2(0,\tau)\theta_3(0,\tau)\theta_4(0,\tau).
\end{equation}
两边均不为零。
\end{proposition}
\begin{proof}
令 \(q=e^{\pi\mathrm i\tau}\)，并规定
\(q^{1/4}=e^{\pi\mathrm i\tau/4}\)。在式\eqref{b3:eq:24-theta1-product}中除以 \(z\) 并令 \(z\to0\)，由于乘积局部一致收敛及 \(\sin z/z\to1\)，得到
\[
 \theta_1'(0,\tau)
 =2q^{1/4}\prod_{n\ge1}(1-q^{2n})^3.
\]
现在将三个 theta 常数的乘积式相乘。共同因子已经给出右边这个式子，剩下的因子是
\[
 \left(\prod_{n\ge1}
    (1+q^{2n})(1-q^{4n-2})\right)^2.
\]
记 \(A=\prod_{n\ge1}(1-q^{2n})\ne0\)。由
\(1+q^{2n}=(1-q^{4n})/(1-q^{2n})\)，括号内的乘积等于
\[
 \frac{\displaystyle
        \prod_{n\ge1}(1-q^{4n})
        \prod_{n\ge1}(1-q^{4n-2})}
      {\displaystyle\prod_{n\ge1}(1-q^{2n})}=1.
\]
分子两列指数 \(4,8,12,\ldots\) 与 \(2,6,10,\ldots\) 正好分割全部正偶数。这里的分组和约分合法，因为各乘积绝对收敛且非零；并没有把两个发散乘积相除。恒等式随之成立，非零性也从这些乘积直接得到。
\end{proof}

这个归一化没有额外的 \(\pi\) 因子，因为本章级数含
\(e^{2\mathrm inz}\)，基本水平准周期为 \(\pi\)。若改用周期为一的空间坐标，空间求导的换元因子也必须随之改变。

当 \(\tau=\mathrm it\)、\(t>0\) 时，\(q=e^{-\pi t}\in(0,1)\)。
三个常数乘积中的全部因子都是正数，因此
\(\theta_2,\theta_3,\theta_4>0\)。尤其是交替级数
\[
 \theta_4=1-2q+2q^4-2q^9+\cdots
\]
始终为正。单看前两项，在 \(q>1/2\) 时容易作出错误判断；乘积同时控制了被截去的无限尾部。对一般复参数，仍有非零性，但不再谈正负。

在 \(q\to0\) 时，级数还给出
\[
 \begin{aligned}
 \theta_2&=2q^{1/4}\bigl(1+q^2+O(q^6)\bigr),\\
 \theta_3&=1+2q+O(q^4),\qquad
 \theta_4=1-2q+O(q^4).
 \end{aligned}
\]
因此它们的乘积为 \(2q^{1/4}(1-3q^2+O(q^4))\)，与导数乘积的展开相符。这个检查能辨认归一化和低次系数，却不能代替对所有 \(\tau\in\mathfrak H\) 成立的恒等式证明。

\subsection{三重积的系数、零点与对数导数}
\label{b3:subsec:240204}

三重积允许比较系数，但要先说清楚比较的是哪个变量。固定
\(w\ne0\) 后，三重积两边关于 \(q\) 在 \(|q|<1\) 全纯，并在
\(q=0\) 都取值一；局部一致收敛保证可以使用幂级数的唯一性。
固定 \(q\) 后，关于 \(w\) 则是穿孔平面上的 Laurent 展开。两种比较都依赖解析收敛。靠近 \(|q|=1\) 时，上面的共同控制不再一致，不能未经新的估计就把 \(q\) 代成单位根。

还有一种需要格外区分的操作：对乘积求对数导数。含有零点的整函数没有全平面上的全纯对数，但在每个不含零点的小圆盘中都有局部对数；不同分支只差常数，所以对数的导数不依赖分支。

\Needspace{8\baselineskip}
\begin{proposition}[theta 的对数导数]\label{b3:prop:24-theta-log-derivative}
设 \(\tau\in\mathfrak H\)、\(q=e^{\pi\mathrm i\tau}\)，令
\(\Lambda=\pi\mathbb Z+\pi\tau\mathbb Z\) 及
\[
 L(z)=\frac{\partial_z\theta_1(z,\tau)}{\theta_1(z,\tau)}
 \qquad(z\in\mathbb C\setminus\Lambda).
\]
则
\begin{equation}\label{b3:eq:24-theta-logarithmic-series}
 L(z)=\cot z+
 4\sin(2z)\sum_{n\ge1}
 \frac{q^{2n}}{1-2q^{2n}\cos(2z)+q^{4n}}.
\end{equation}
右边的级数在 \(\mathbb C\setminus\Lambda\) 的紧集上绝对一致收敛，并可作任意固定阶数的逐项微分。在更小的范围
\(|\operatorname{Im}z|<\pi\operatorname{Im}\tau\)、\(z\notin\pi\mathbb Z\) 内，还有
\begin{equation}\label{b3:eq:24-theta-log-fourier}
 L(z)=\cot z+
 4\sum_{k\ge1}\frac{q^{2k}}{1-q^{2k}}\sin(2kz).
\end{equation}
函数 \(-L'\) 是以 \(\pi,\pi\tau\) 为周期的亚纯函数，其全部极点是
\(\Lambda\)，在每个 \(\lambda\in\Lambda\) 处的主部都是
\((z-\lambda)^{-2}\)。
\end{proposition}
\begin{proof}
在避开 \(\Lambda\) 的紧集上，theta 乘积的有限头部非零，尾因子一致接近一。把紧集稍微放大后，尾因子的偏差仍由几何级数控制。因此\cref{b3:prop:20-product-log-derivative}的条件成立：可对局部对数求导，并可继续逐项微分。由
\[
 \frac{\mathrm d}{\mathrm dz}
   \log(1-2q^{2n}\cos2z+q^{4n})
 =\frac{4q^{2n}\sin2z}
        {1-2q^{2n}\cos2z+q^{4n}},
\]
以及首因子 \(\sin z\) 的对数导数，得到第一式。这里的公式只要求避开零点，不要求 \(z\) 留在一个水平带内。

第二式需要额外的几何级数展开。若
\(|\operatorname{Im}z|<-\log|q|=\pi\operatorname{Im}\tau\)，则对每个 \(n\ge1\)，
\(|q^{2n}e^{\pm2\mathrm iz}|<1\)。将两个指数因子的对数导数分别展开后相加，得到
\[
 \frac{-2\mathrm i q^{2n}e^{2\mathrm iz}}
           {1-q^{2n}e^{2\mathrm iz}}
 +\frac{2\mathrm i q^{2n}e^{-2\mathrm iz}}
           {1-q^{2n}e^{-2\mathrm iz}}
 =4\sum_{k\ge1}q^{2nk}\sin(2kz).
\]
在闭子带 \(|\operatorname{Im}z|\le R<-\log|q|\) 上，相应双重级数的绝对值和不超过常数乘以
\[
 \sum_{k\ge1}
   \frac{|q|^{2k}e^{2kR}}{1-|q|^{2k}}<\infty.
\]
故可交换 \(n,k\) 的求和顺序，先对 \(n\) 求和，就得到式\eqref{b3:eq:24-theta-log-fourier}。

由上一节的准周期公式，
\[
 \theta_1(z+\pi,\tau)=-\theta_1(z,\tau),\qquad
 \theta_1(z+\pi\tau,\tau)
 =-q^{-1}e^{-2\mathrm iz}\theta_1(z,\tau).
\]
取对数导数即得 \(L(z+\pi)=L(z)\) 和
\(L(z+\pi\tau)=L(z)-2\mathrm i\)，再求导便消去加法常数。
另一方面，\cref{b3:prop:24-theta-zeros}保证每个格点均为单零点。
在 \(\lambda\) 附近可写
\(\theta_1(z,\tau)=(z-\lambda)h(z)\)，其中 \(h\) 全纯且非零。于是
\[
 L(z)=\frac1{z-\lambda}+\frac{h'(z)}{h(z)},\qquad
 -L'(z)=\frac1{(z-\lambda)^2}
          -\left(\frac{h'}h\right)'(z).
\]
第二项全纯，从而得到全部极点和主部。
\end{proof}

例如，在原点附近对式\eqref{b3:eq:24-theta-log-fourier}求导，并使用
\(\csc^2z=z^{-2}+\tfrac13+O(z^2)\)，得到
\begin{equation}\label{b3:eq:24-theta-local-part}
 -L'(z)=\frac1{z^2}+\frac13
       -8\sum_{k\ge1}\frac{kq^{2k}}{1-q^{2k}}+O(z^2).
\end{equation}
这里极点的系数一已经由单零点确定，常数项则由乘积归一化决定。后文与同一格的 \(\wp\) 函数比较时，二者之差只能剩下常数；不能只因二者都双周期就认定它们相等。

参数方向的对数微分也要有共同控制。例如
\[
 B(q)=\prod_{n\ge1}(1-q^{2n})\qquad(|q|<1)
\]
处处非零且全纯，并满足
\[
 q\frac{B'(q)}{B(q)}
 =-\sum_{n\ge1}\frac{2nq^{2n}}{1-q^{2n}}.
\]
在 \(|q|\le r<1\) 上，右边由
\(\sum_{n\ge1}2nr^{2n}/(1-r^{2n})<\infty\) 控制。
其他 theta 常数乘积的参数导数有同样的几何衰减。若出现
\(q^{1/4}\)，则在 \(\tau\) 平面使用
\[
 q\partial_q=\frac1{\pi\mathrm i}\partial_\tau,
 \qquad q\partial_q\log(q^{1/4})=\frac14.
\]
后一式是局部对数的导数陈述，不要求在整个穿孔 \(q\) 圆盘选定单值的 \(\log q\)。这些估计将使后文从乘积转到约数和的逐项运算有明确的适用范围。
